\section{Introduction}\label{sec:introduction}
Vision-language-action (VLA) models connect visual observations and language instructions to robot commands. Their pretrained representations offer a starting point for manipulation, but their action interface depends on the embodiment used during training. When a new robot requires different action dimensions or control conventions, adapting the model can involve replacing or reinitializing its action head or expert. In this setting, useful visual and language representations may be available before the policy can reliably execute even a basic movement.

Learning a complete instruction directly from binary success is difficult when the initial policy rarely succeeds. Consider ``put the apple on the plate'': the robot must first approach the apple, grasp it, lift it, carry it, and release it correctly. A failed rollout gives the same final outcome whether the robot never reaches the object or drops it just before placement. Training each component separately can make learning easier, but those components must subsequently work together from a normal task start under one instruction.

Our approach follows this structure. We first use dense rewards to learn individual skills, giving the policy feedback before it can complete them. We then transition to binary skill-success rewards and self-imitation of successful experience. Finally, we train the complete task with staged rewards so that earlier skills and later task progress contribute to learning. The aim is to adapt VLA-based control to a new embodiment through a practical training sequence, rather than to make a particular robot architecture the subject of the method.

We investigate embodiment adaptation using SmolVLA and a cable-driven parallel robot (CDPR), whose gripper is positioned by coordinated cable actuation. This platform serves as a demanding example of a robot with a different mechanical and control structure from conventional serial arms. Its role is to test adaptation under a substantial embodiment change; the approach does not rely on a cable-specific task objective. The evaluated implementation learns a new residual action module while retaining the pretrained SmolVLA action expert and representations (Sec.~\ref{sec:method}).

We evaluate two aspects of this approach. Our OpenVLA-OFT experiments on LIBERO examine reward construction: phase transitions and newly achieved within-phase progress provide intermediate feedback while preserving the terminal verifier reward. Our CDPR experiments examine embodiment adaptation and full-task composition using discrete milestones and stage-specific advantages. Both expose ordered manipulation progress to learning and evaluate task completion. The reward experiments ask whether intermediate feedback remains useful with a trained action policy; the adaptation experiments ask whether the learning sequence yields complete-task control on the target embodiment.

The paper contributes a dense-to-binary skill-learning schedule with self-imitation and staged composition; phase-aware intermediate feedback with dense-progress and milestone implementations; and evaluations of reward construction and complete-task execution. The complete schedule is instantiated on the CDPR, while LIBERO evaluates the reward design with an established action interface. This distinction defines what each experiment can establish about the approach.
